\section{System Architecture and Data}
\label{sec:system}

\subsection{Architecture}

HospitAI is divided into a web application and an offline analytics part. The web side has a frontend built with Next.js and a backend built with FastAPI. The backend follows a strict layered structure: API routers call a service layer, the service layer calls a repository layer, and only the repository layer talks to the database. This separation keeps the database queries in one place and makes the code easier to maintain. All persistent data lives in a single PostgreSQL database. Figure~\ref{fig:arch} shows the main components.

The two heavy machine-learning steps, clustering and sequence mining, do not run inside a web request. They are separate Python scripts that we run offline, because they are slow and only need to run when the data changes. Each script reads from the database, computes its result and writes it back into dedicated tables. The web application then only reads these tables, which keeps the user interface responsive. A single LLM, accessed through the OpenAI API, is used for every text-generation task in the platform.

\begin{figure}[t]
\centering
\begin{tikzpicture}[node distance=6mm and 9mm]
  \node[box, fill=blue!8] (fe) {Frontend (Next.js + Redux)};
  \node[box, below=of fe] (api) {Backend API (FastAPI routers)};
  \node[box, below=of api] (svc) {Service layer};
  \node[box, below=of svc] (repo) {Repository layer};
  \node[db, below=9mm of repo] (db) {PostgreSQL};
  \node[step, right=16mm of svc, text width=26mm] (ml) {Offline ML scripts:\\K-Prototypes\\PrefixSpan};
  \node[llm, left=12mm of svc, text width=20mm] (llm) {LLM\\(OpenAI GPT)};

  \draw[flow] (fe) -- (api);
  \draw[flow] (api) -- (svc);
  \draw[flow] (svc) -- (repo);
  \draw[flow] (repo) -- (db);
  \draw[flow] (svc) -- (llm);
  \draw[{Stealth[length=2mm]}-{Stealth[length=2mm]}, draw=black!70] (ml) |- (db)
       node[pos=0.28,right,font=\footnotesize]{read / write};
\end{tikzpicture}
\caption{High-level architecture. The web application reads precomputed results, while the clustering and mining scripts run offline and write their output into the database.}
\label{fig:arch}
\end{figure}

\subsection{Data}

For development and experiments we use a synthetic dataset generated with Synthea, an open-source patient simulator. Synthetic data is a sensible choice for a student project, because it removes privacy and legal barriers while keeping a realistic clinical structure. The dataset contains about 22{,}284 patients, with the usual clinical tables: encounters, conditions, medications, procedures, observations and care plans. A loading script reduces the data to a single hospital and a fixed set of providers, cleans the dates, and computes helper fields such as patient age and treatment duration.

The dataset is large enough to make the analytics non-trivial. The event tables used for sequence mining hold roughly 727{,}000 conditions, 913{,}000 medications, 3.2 million procedures and 70{,}000 care plans, which together give close to five million clinical events. After cleaning, the vocabulary of distinct event items is around 1{,}170.

The same database stores both the raw clinical tables and the analytics output. Clustering writes one row per patient with the cluster label, a flag for outliers, a two-dimensional projection for visualisation, and a cached set of features in a JSON column. Sequence mining writes the discovered patterns into one table and the patient-to-pattern membership into a link table. The interpretation step stores its generated text in JSON columns next to the cluster and pattern rows, so that each summary is computed only once and then reused.
